\begin{definition}\label{def:bola} 
    \textbf{Bola de radio $\delta$ y centro $\mathbf{x}_0$.} 
    Dados $\mathbf{x}_0 \in \Rn{n}$, $\delta > 0$ y $d$ una función distancia en $\Rn{n}$.  Se define la \emph{bola abierta de radio $\delta$ y centro $\mathbf{x}_0$}, que denotamos con $B_{\delta}(\mathbf{x}_0)$, como el conjunto
    \[
     B_{\delta}(\mathbf{x}_0) = \left\lbrace \mathbf{x} \in \Rn{n} : d(\mathbf{x},\mathbf{x}_0) < \delta \right\rbrace.
    \]
    El conjunto
    \[
     B^*_{\delta}(\mathbf{x}_0) = B_{\delta}(\mathbf{x}_0) \smallsetminus \{ \mathbf{x}_0 \}
    \]
    se llama \emph{bola reducida de radio $\delta$ y centro $\mathbf{x}_0$}.
\end{definition}

\begin{example}
    Para ilustrar lo mencionado, podemos tomar el caso de dos bolas en $\Rn{2}$ centradas en el origen, de ``radio'' $\delta$ y con distintas nociones de distancia. La primera con distancia inducida por la norma 2 y la segunda con distancia inducida por la norma 1. Se deja al lector deducir anal\'iticamente los gr\'aficos mostrados a continuaci\'on. 
\input{figs/bola.tex}
\end{example}



  \begin{definition} \textbf{Entorno de un punto.} \label{def:entorno}
    Sean $N  \subseteq   \Rn{n}$ y $\vect{x}_0 \in \Rn{n}$. Decimos que $N$ es \emph{entorno} de $\vect{x}_0$ si
    \[
     \exists \delta > 0 \; \colon \; B_{\delta}(\vect{x}_0)  \subseteq  N.
    \]
   \end{definition}

 \emph{Notaci\'on:} De ahora en m\'as, denotaremos a la norma 2 simplemente como $\|\cdot\|$, ya que va a ser la \'unica norma utilizada a lo largo del documento.
   